\chapter{Convergence proof for the allocation CRDT}\label{sec:appendix_proof}

This appendix proves that the topic allocation CRDT described in
section~\ref{sec:session_crdt} converges to a collision-free, non-divergent
state in finite time, under the assumptions A1--A9 restated below.
The proof applies both to the formal operators in
appendix~\ref{sec:appendix_core_tla} (\texttt{AllocateTopic},
\texttt{AcceptGossip}, \texttt{Converged}) and to the C implementation in
\texttt{cy.c} (\texttt{topic\_allocate}, \texttt{on\_gossip\_known\_topic},
\texttt{on\_gossip\_unknown\_topic}) under the stated assumptions.

\section{Model}\label{sec:appendix_proof_model}

Let:

\begin{itemize}
    \item $H$ be the finite set of active topic hashes, with $|H| = N$.
    \item $U$ be the finite set of auto-assigned subject-IDs, with $|U| = M$.
    \item $S : H \times \mathbb{N} \to U$ be the deterministic subject-ID
    mapping function.
    \item A topic record be a triple $(h, e, a)$ where $h \in H$,
    $e \in \mathbb{N}$ is the eviction counter, and $a \in \mathbb{N}$ is
    the age.
    \item $\ell(a) = \lfloor \log_2 a \rfloor$ with $\ell(0) = -1$ be the
    log-age function.
\end{itemize}

Two records are in \emph{collision} when they have different hashes but
identical subject-IDs. Two records are in \emph{divergence} when they share
a hash but carry different eviction counters. A state is \emph{converged}
when it contains no collisions and no divergences.

\section{Assumptions}\label{sec:appendix_proof_assumptions}

Some of the following are structural properties of the mapping function and
are satisfied by construction; others are deployment constraints that must
be upheld by compliant implementations and system integrators.

\begin{description}
    \item[A1. Common deterministic mapping.] All nodes use the same
    deterministic $S(h, e)$ and the same mapping parameters -- in particular,
    the same modulus.

    \item[A2. Hash identity.] Active topic names are hash-unique
    (equivalently: equal hash means the same logical topic).
    Hash collision probability is negligible and is ignored by the protocol.

    \item[A3. Bounded escape property.] There exists $K \in \mathbb{N}_{>0}$
    such that
    \[
        \forall h \in H,\ \forall x \in \mathbb{N},\ \forall B \subseteq U,\
        |B| < K,\ \exists d \in \{0, \dots, K-1\} :
        S(h,\ x + d) \notin B.
    \]
    In prose: from any current eviction value, a topic can escape any
    blocked subject-ID set smaller than $K$ within at most $K$ probing
    increments.

    \item[A4. Capacity bound.] $N \le K$.

    \item[A5. Gossip fairness.] Every live allocation entry is eventually
    received and processed by every relevant node, infinitely often.

    \item[A6. Progress fairness.] Each node keeps executing the protocol
    loop -- no permanent local halt while the node is considered active.

    \item[A7. Resource sufficiency.] Memory, stack, and scheduler resources
    are sufficient for processing each local allocation cascade to
    completion.

    \item[A8. Eviction counter safety.] Eviction counters do not overflow
    in reachable executions.

    \item[A9. Stability window.] There exists a repair bound
    $T_{\mathrm{repair}}$ and time windows of length greater than
    $T_{\mathrm{repair}}$ during which the active topic set $H$ is
    unchanged, the mapping parameters are unchanged, and pairwise
    arbitration outcomes do not flip. Equivalently: arbitration-priority
    flips driven by $\lfloor \log_2 a \rfloor$ increments occur much less
    frequently than the gossip-plus-repair dynamics.
\end{description}

\section{Theorem 1: Convergence in one stability window}

Under assumptions A1--A9, any execution restricted to a single stability
window converges in finite time to a collision-free, non-divergent state.

\subsection*{Lemma 1: Local allocation termination under bounded escape}

Fix a node and a moving topic hash $h$ with current eviction $x$. Let $B$
be the set of subject-IDs currently occupied by other local topics.
If $|B| < K$, then by A3 there exists $d < K$ such that
$S(h, x+d) \notin B$.
\texttt{AllocateTopic} increments the moving topic's eviction counter by
one each time it loses arbitration on its current candidate subject-ID.
Therefore within at most $d$ losses the topic reaches a subject-ID not in
$B$, and the insertion terminates.

Local allocation terminates whenever $|B| < K$. By A4, $|B| \le N - 1 < K$,
so allocation always terminates.

\subsection*{Lemma 2: Divergence repair converges for each hash}

Fix a hash $h$. \texttt{AcceptGossip} on the known-hash path compares the
local and remote replicas using divergence arbitration and either keeps the
local evictions (merging the age) or adopts the remote evictions and
re-allocates.

Under A5 and within a stable window (A9), every node repeatedly observes
the same finite set of contenders for $h$ and eventually agrees on the same
winning replica -- the same eviction value and a consistent merged
log-age. Hence all divergences for $h$ disappear in finite time. Applying
this reasoning independently to each $h \in H$, all divergences disappear
in finite time.

\subsection*{Lemma 3: Collision repair converges under stable priorities}

In a stability window (A9), collision arbitration induces a fixed strict
total order on the active hashes:
\[
    h_1 \succ h_2 \succ \dots \succ h_N.
\]
When $h_i$ collides with any $h_j$ where $j < i$, $h_i$ loses and is
evicted; $h_j$ is unchanged. Higher-priority topics are therefore never
displaced by lower-priority topics.

Induction on the priority rank:
\begin{itemize}
    \item \textbf{Base $i = 1$.} $h_1$ cannot be displaced by anyone, so
    once allocated it remains fixed.
    \item \textbf{Inductive step.} Assume $h_1, \dots, h_{i-1}$ are fixed
    at a subject-ID set $B_i$ with $|B_i| = i - 1 < K$ (by A4). By A3,
    from the current eviction value $x_i$ topic $h_i$ has some $d_i < K$
    with $S(h_i,\ x_i + d_i) \notin B_i$. Each loss against $B_i$
    increments the eviction counter by one, so after at most $d_i$ losses
    $h_i$ reaches a subject-ID outside $B_i$ and cannot be displaced by
    any lower-priority topic.
\end{itemize}

Every $h_i$ stabilizes in finite time, so all collisions disappear.

\subsection*{Conclusion}

By Lemma~2 divergences vanish in finite time. By Lemma~3 collisions vanish
in finite time. The state therefore becomes converged -- no collisions, no
divergences -- in finite time.

\section{Theorem 2: Convergence with occasional priority flips}

Assume A1--A8 and that executions consist of stability windows as in A9,
each longer than $T_{\mathrm{repair}}$. Then the network converges within
each window; non-converged periods are transient and confined to window
transitions (including the effects of in-flight stale gossip).

\subsection*{Proof sketch}

Apply Theorem~1 independently to each stability window. Window transitions
may temporarily reintroduce conflicts, but by construction the next stable
window is long enough to repair them before the following flip
($\text{window length} > T_{\mathrm{repair}}$).

\section{Instantiating A3/A4 for common probing laws}

\paragraph{Linear probing} $S(h, e) = h + e \bmod M$ satisfies A3 with
$K = M$. This is the instantiation used in the TLA\textsuperscript{+}
model of appendix~\ref{sec:appendix_core_tla} for legibility.

\paragraph{Quadratic probing} This is the instantiation used by the actual
Cyphal v1.1 implementation (section~\ref{sec:session_subject_id_mapping}).
The protocol relies on the deployment bound that the number of active
non-pinned topics does not exceed the collision-free probing capacity;
practically this is the half-ring bound $K \approx \lfloor M / 2 \rfloor + 1$
and consequently $N \le \lfloor M / 2 \rfloor$.
